\section{Reproducing the numbers}
\label{app:repro}

Every number quoted in this paper is generated, not typed. The prose reads its
values from a single generated macro file, so a figure and the sentence describing
it cannot drift apart: if an experiment is re-run and its result changes, the text
changes with it or the build fails.

\subsection{Layout}

Experiments live in \texttt{paper\pslash tools\pslash exp\pslash}, one file per
claim, each writing a JSON or CSV file to \texttt{paper\pslash data\pslash} and its
figures to \texttt{paper\pslash figures\pslash}. A collector reads every data file and emits
\texttt{numbers.tex}, a list of macros; the prose uses only those macros for
quantities. The mapping from claim to command is:

\begin{center}\footnotesize
\begin{tabular}{@{}lll@{}}
  \toprule
  Claim & Script & Data \\
  \midrule
  Fringe law (Thm.~\ref{thm:fringe})     & \texttt{fringe.mjs}     & \texttt{fringe.json} \\
  Fringe law, drawn (Fig.~\ref{fig:visibility}) & \texttt{fringelaw.mjs} &---\\
  Law across pairs (Fig.~\ref{fig:lawatlas}) & \texttt{lawatlas.mjs} & \texttt{lawatlas.json} \\
  Eikonal divide (\S\ref{sec:eikonal})   & \texttt{gradient.mjs}   & \texttt{gradient.json} \\
  Cost (Table~\ref{tab:cost})            & \texttt{cost.mjs}       & \texttt{cost.json} \\
  Fidelity (Table~\ref{tab:fidelity})    & \texttt{fidelity.mjs}   & \texttt{fidelity.json} \\
  Ablation (Table~\ref{tab:ablation})    & \texttt{gpu\pslash ablation.mjs} & \texttt{ablation.json} \\
  Interval width (\S\ref{sec:limits})    & \texttt{saturation.mjs}   & \texttt{saturation.json} \\
  Degenerate band (\S\ref{sec:limits})   & \texttt{degenerate.mjs} & \texttt{degenerate.json} \\
  Bounds, drawn (Figs.~\ref{fig:bounds},~\ref{fig:convex}) & \texttt{math.mjs} & \texttt{residual.csv} \\
  Interval verified (\S\ref{sec:window}) & \texttt{math.mjs}       & \texttt{window.json} \\
  Artifacts (Figs.~\ref{fig:holes}--\ref{fig:anchor}) & \texttt{artifacts.mjs} & \texttt{artifacts.json} \\
  Selection (\S\ref{sec:selection})      & \texttt{convergents.mjs} & \texttt{convergents.json} \\
  Selection, drawn (Figs.~\ref{fig:selection},~\ref{fig:pitch3}) & \texttt{selectionfigs.mjs} & --- \\
  Defects (\S\ref{sec:defects})          & \texttt{defects.mjs}    & \texttt{defects.json} \\
  Defects, drawn (Fig.~\ref{fig:defects}) & \texttt{defectfigs.mjs} & --- \\
  Contouring (\S\ref{sec:beyond})        & \texttt{contour.mjs}    & \texttt{contour.json} \\
  Field expressions (Table~\ref{tab:compile}) & \texttt{gpu\pslash run.mjs} & \texttt{gpu.json} \\
  Stroke floor (\S\ref{sec:zoom})        & \texttt{zoom.mjs}       & \texttt{zoom.json} \\
  Explicit geometry (\S\ref{sec:vector}) & \texttt{vector.mjs}     & \texttt{vector.json} \\
  Catalog plate (Fig.~\ref{fig:plate})   & \texttt{plate.mjs}      &---\\
  Two traditions (Fig.~\ref{fig:traditions}) & \texttt{traditions.mjs} &---\\
  Teaser (Fig.~\ref{fig:teaser})         & \texttt{teaser.mjs}     &---\\
  \bottomrule
\end{tabular}
\end{center}

\subsection{The solvers under test are the shipped sources}

The comparisons in \S\ref{sec:results} run four solvers, and none is a
reimplementation. \sweep{} and \loose{} are historical versions, kept
verbatim under \texttt{paper\pslash tools\pslash legacy\pslash}; \ours{} is read
live from \texttt{src\pslash gpu\pslash inverseCpu.ts}, the file the shader's CPU
twin actually uses, so the paper always measures what ships; and
\textsc{Reference} is the exhaustive scan of \S\ref{sec:protocol}. All are
loaded by rewriting their source at import time to wrap
\texttt{shapeRadius} in a counter, so ``metric evaluations per pixel'' is measured
in the shipping code rather than estimated beside it. An ablation is expressed the
same way, as a literal find-and-replace against the shipping source; a patch whose
pattern fails to match raises rather than silently doing nothing, which is what
keeps the ablation table honest as the source moves underneath it.

Because the counter wraps the metric rather than the loop, the unit compares across
solvers that search differently: in every solver, one candidate index costs one
\texttt{shapeRadius} call and one Newton step a \texttt{shapeRadius} plus a
gradient.

Table~\ref{tab:compile} is arranged the same way. The compiled column is produced
by calling the shipping emitter, \texttt{src\pslash fields\pslash emit.ts}, and
pasting its output into the timing kernel, so it is the code the tool runs and not
a hand transcription of it. The interpreter it is compared against is the shader we
replaced, kept verbatim under \texttt{paper\pslash tools\pslash legacy\pslash}. Both
are checked against \texttt{src\pslash fields\pslash evalExpr.ts} at the same points
before either is timed.

\subsection{Figures are renders, not screenshots}

No figure in this paper is a capture of the running application. The pattern figures
come from a CPU rasterizer that mirrors the shader's \texttt{composite.ts} term for
term (same world transform, same $1.15$-pixel stroke floor, same $0.7$-pixel
antialiasing band, same accept/reject guards, same alpha composite), so a figure
can be rendered with a solver that no longer exists in the shader, and per-pixel
work can be counted while it is drawn. Where a figure needs the walking families it
calls the shipped solver through the experiments' own instrumented loader.

This costs orders of magnitude more time than a screenshot and buys exact
reproducibility, antialiasing included. It matters for one claim: the sub-pixel
offset measured in \S\ref{sec:beyond} would otherwise be confounded by an unknown
display pipeline.

\subsection{The clips are renders too}

The moving figures in the supplemental video are made the same way and for the
same reason. Each frame is drawn by the shipped renderer at an explicit
timestamp and an explicit size (frame $n$ at $t_0+n/\mathrm{fps}$, never at
whatever moment the browser's animation callback happened to reach) and handed
to the encoder in order, so the rate a clip claims is the rate it has. Nothing
is captured from a screen. A clip is therefore a function of the range it names
rather than of the machine that made it, and re-recording the same range from
the same scene file gives the same frames. This is the same property
\S\ref{sec:time} relies on: the renderer holds no state, so a frame is
determined by the clock and by nothing else.

\subsection{Running it}

The experiments need only a recent Node with TypeScript stripping enabled, since
they import the shader's twin directly. Run the CPU twin's unit tests first, because
everything else depends on them; then any experiment; then the collector and the
document build. Total runtime is dominated by \texttt{fidelity.mjs} and
\texttt{saturation.mjs}, which sweep $\fidSettings$ and $\satSettings$ settings
respectively against an exhaustive reference.

\begin{center}\footnotesize
\begin{tabular}{@{}l@{}}
  \texttt{node -{}-experimental-strip-types src\pslash gpu\pslash inverseCpu.test.ts} \\
  \texttt{node paper\pslash tools\pslash exp\pslash contour.mjs} \\
  \texttt{node paper\pslash tools\pslash numbers.mjs} \\
\end{tabular}
\end{center}
